\documentclass[../../../include/open-logic-section]{subfiles}

\begin{document}

\olfileid{sth}{story}{extensionality}
\olsection{في الامتدادية}

أول أصول هذا الباب أن هوية المجموعة تتعين بما ينتمي إليها من
!!{element}. وتقرير ذلك على وجه الدقة:

\begin{axiom}[الامتدادية]
إذا اشتركت المجموعتان~$A$ و$B$ في !!{element} نفسها، كانت~$A$ و$B$
مجموعة واحدة.
\[
  \lforall[A][\lforall[B][(\lforall[x][(x \in A \liff x \in B)] \lif
  \eq[A][B])]]
\]
\end{axiom}

وقد جرى افتراض هذا الأصل في~\olref[sfr][][]{part} كله، غير أن إعادته
هنا واجبة. فمسلّمة الامتدادية تقول إن المجموعة لا يحددها إلا ما ينتمي
إليها من !!{element}. ومن هنا تفارق المجموعات \emph{المفاهيم}، إذ قد
تندرج الأشياء عينها تحت مفاهيم كثيرة متباينة.

ولِمَ نقبل هذا المبدأ؟ لأن ردّ الامتدادية عدول عما يصح أن يسمى
\emph{نظرية المجموعات}؛ فما هذه النظرية إلا نظرية التجمعات الامتدادية.

وإنما موضع العسر في~\olref[sth][][]{part} أن نضع أصولًا تعيّن
\emph{المجموعات الموجودة}. وسنرى أن الجواب الذي يبدو «بديهيًا» حقًا
باطل ببرهان.

\end{document}



